\begingroup
\renewcommand{\arraystretch}{1.45}
\fancyhead[L]{\small ARITHMETIC–GEOMETRIC MEAN.}
\Needspace{28\baselineskip}
\AGMArticle{25}
[From the functional equations for $P,\allowbreak Q,\allowbreak R,\allowbreak S$ that hold when the second argument $\eta$ is changed into its reciprocal, when its sign is changed and when it is multiplied by the first argument $y$, as well as from the known values of the functions for $\eta=1$, or also from the partial differential equation established in Article 21 for the general functions, it follows that, if $P,\allowbreak Q,\allowbreak R,\allowbreak S$ can be expanded in series in increasing integer powers of $y^{\frac14}$ and $\eta^{\frac12}$, these must be
\Needspace{30\baselineskip}
\[
\begin{aligned}
P(y,\eta)&=1+y(\eta+\eta^{-1})\\[4pt]&\quad{}+y^4(\eta^2+\eta^{-2})\\[4pt]&\quad{}+y^9(\eta^3+\eta^{-3})\\[4pt]&\quad{}+y^{16}(\eta^4+\eta^{-4})+\,+\,.\\[10pt]
Q(y,\eta)&=1-y(\eta+\eta^{-1})\\[4pt]&\quad{}+y^4(\eta^2+\eta^{-2})\\[4pt]&\quad{}-y^9(\eta^3+\eta^{-3})\\[4pt]&\quad{}+y^{16}(\eta^4+\eta^{-4})-\,+\,.\\[10pt]
R(y,\eta)&=y^{\frac14}(\eta^{\frac12}+\eta^{-\frac12})+y^{\frac94}(\eta^{\frac32}+\eta^{-\frac32})\\[4pt]&\quad{}+y^{\frac{25}4}(\eta^{\frac52}+\eta^{-\frac52})\\[4pt]&\quad{}+y^{\frac{49}4}(\eta^{\frac72}+\eta^{-\frac72})+\,+\,.\\[10pt]
S(y,\eta)&=y^{\frac14}(\eta^{\frac12}-\eta^{-\frac12})-y^{\frac94}(\eta^{\frac32}-\eta^{-\frac32})\\[4pt]&\quad{}+y^{\frac{25}4}(\eta^{\frac52}-\eta^{-\frac52})\\[4pt]&\quad{}-y^{\frac{49}4}(\eta^{\frac72}-\eta^{-\frac72})+\,-\,.
\end{aligned}
\]

\Needspace{24\baselineskip}
The statement that the series $P,\allowbreak Q$ multiplied by the limit $\surd  H$ represent the quantities $\surd  A,\allowbreak \surd  B$, to which the algorithm used by \textsc{Gauss} and reproduced at the end of Article 18 refers, is given as a special theorem in the manuscript remains, and at the same time it is noted that]
\SourceLeaf{409}

\Needspace{24\baselineskip}
\[
\begin{aligned}
\operatorname{tang}\tfrac12 U
 &=\dfrac{\begin{gathered}y^{\frac18}\sin\dfrac12u-y^{\frac98}\sin\dfrac32u\\[3pt]+y^{\frac{25}8}\sin\dfrac52u-y^{\frac{49}8}\sin\dfrac72u+\mathord{.}\,\mathord{.}\end{gathered}}{\begin{gathered}y^{\frac18}\cos\dfrac12u+y^{\frac98}\cos\dfrac32u\\[3pt]+y^{\frac{25}8}\cos\dfrac52u+y^{\frac{49}8}\cos\dfrac72u+\mathord{.}\,\mathord{.}\end{gathered}}\\[14pt]
\operatorname{tang} U'
 &=\dfrac{\begin{gathered}y^{\frac14}\sin u-y^{\frac94}\sin3u\\[3pt]+y^{\frac{25}4}\sin5u-y^{\frac{49}4}\sin7u+\mathord{.}\,\mathord{.}\end{gathered}}{\begin{gathered}y^{\frac14}\cos u+y^{\frac94}\cos3u\\[3pt]+y^{\frac{25}4}\cos5u+y^{\frac{49}4}\cos7u+\mathord{.}\,\mathord{.}\end{gathered}}\\[14pt]
\operatorname{tang}2U''
 &=\dfrac{\begin{gathered}y^{\frac12}\sin2u-y^{\frac92}\sin6u\\[3pt]+y^{\frac{25}2}\sin10u-y^{\frac{49}2}\sin14u+\mathord{.}\,\mathord{.}\end{gathered}}{\begin{gathered}y^{\frac12}\cos2u+y^{\frac92}\cos6u\\[3pt]+y^{\frac{25}2}\cos10u+y^{\frac{49}2}\cos14u+\mathord{.}\,\mathord{.}\end{gathered}}
\end{aligned}
\]
[results. For the theorem that the series $P$, $Q$ satisfy the functional equations that determine the algorithm discussed, and that were collected above in Article 22, \textsc{Gauss} probably also gave another proof, besides the proof based on transforming those series into infinite products and contained in the paper ‘Hundred Theorems on the New Transcendent’ that follows below, such as can easily be obtained from the indications given above in Article 16.]

\Needspace{24\baselineskip}
\AGMArticle{26}
[If, departing from the notation followed in the preceding articles, we denote the first derivatives of the functions $P(y,\allowbreak \eta)$, $Q(y,\allowbreak \eta)$, $R(y,\allowbreak \eta)$, $S(y,\allowbreak \eta)$ with respect to the quantity $\log\eta$ as independent variable by $P'$, $Q'$, $R'$, $S'$ and the second derivatives with respect to the same quantity by $P''$, $Q''$, $R''$, $S''$, and further the first derivatives of the functions $py$, $qy$, $ry$ with respect to $\log y$ as independent variable by $p'$, $q'$, $r'$, then the differential equations found for $a$, $b$, $c$, $\alpha$, $\beta$, $\gamma$, $\delta$ give:
\Needspace{12\baselineskip}
\[
\begin{aligned}\dfrac{qr'-rq'}{qrp^4}&=\dfrac{pr'-rp'}{rpq^4}
 \\[7pt]&=\dfrac{qp'-pq'}{pqr^4}\\[7pt]&=\tfrac14\end{aligned}
\]
\Needspace{20\baselineskip}
\[
\begin{aligned}\dfrac{PS'-SP'}{ppQR}&=\dfrac{QS'-SQ'}{qqPR}
 \\[7pt]&=\dfrac{RS'-SR'}{rrPQ}\\[7pt]&=\dfrac{QR'-RQ'}{ppPS}
 \\[7pt]&=\dfrac{PR'-RP'}{qqQS}\\[7pt]&=\dfrac{QP'-PQ'}{rrRS}\\[7pt]&=1\end{aligned}
\]
\Needspace{12\baselineskip}
\[
\begin{aligned}-4\dfrac{Q'}Q&=(ry)^2\mathbin{.}\dfrac{S(yy,\eta\eta)}{Q(yy,\eta\eta)}\\[6pt]&\quad{}+
(ryy)^2\mathbin{.}\dfrac{S(y^4,\eta^4)}{Q(y^4,\eta^4)}\\[6pt]&\quad{}+
(ry^4)^2\mathbin{.}\dfrac{S(y^8,\eta^8)}{Q(y^8,\eta^8)}+\mathord{.}\,\mathord{.}\end{aligned}
\]
\Needspace{22\baselineskip}
\[
\begin{aligned}
\dfrac{P'P'-PP''}{PP}+\tfrac14pprr\dfrac{RR}{PP}
 &=\dfrac{Q'Q'-QQ''}{QQ}-\tfrac14pprr\dfrac{SS}{QQ}
\\[8pt]&=\dfrac{R'R'-RR''}{RR}+\tfrac14pprr\dfrac{PP}{RR}
 \\[8pt]&=\dfrac{S'S'-SS''}{SS}-\tfrac14pprr\dfrac{QQ}{SS}
\\[8pt]&=-\dfrac{q'}q\\[8pt]&=\tfrac18\{(ry)^4+2(ryy)^4+4(ry^4)^4+8(ry^8)^4+\mathord{.}\,\mathord{.}\}
\end{aligned}
\]

\SourceLeaf{410}

\Needspace{8\baselineskip}
\[
\dfrac{\mathrm d\mathrm dP}{(\mathrm d\log\eta)^2}
 =\dfrac{\mathrm dP}{\mathrm d\log y},
\]
each of the functions $P$, $Q$, $R$, $S$ satisfying this last equation. The preceding chain of equalities between the first and second derivatives of the functions with respect to the quantity $\log\eta$ is found, as far as it concerns $P$ and $S$, in \textsc{Gauss}'s manuscript remains at a place separate from the other investigations of these functions. There $P$, $S$, $p$, $r$ are defined by their series in a form somewhat more convenient for this particular development than the one used here, \Needspace{14\baselineskip}
and the text then reads,] \textit{this gives}
\[
\begin{aligned}
P'P'-PP''&=-\dfrac{\mathrm dp(yy)}{\mathrm d\log y}\mathbin{.}P(yy,\eta\eta)\\[5pt]&\quad{}-
\dfrac{\mathrm dr(yy)}{\mathrm d\log y}\mathbin{.}R(yy,\eta\eta)\\[9pt]
S'S'-SS''&=+\dfrac{\mathrm dr(yy)}{\mathrm d\log y}\mathbin{.}P(yy,\eta\eta)\\[5pt]&\quad{}-
\dfrac{\mathrm dp(yy)}{\mathrm d\log y}\mathbin{.}R(yy,\eta\eta)
\end{aligned}
\]
\Needspace{8\baselineskip}
\[
\begin{aligned}
PP&=+p(yy)\mathbin{.}P(yy,\eta\eta)+r(yy)\mathbin{.}R(yy,\eta\eta)\\
SS&=-r(yy)\mathbin{.}P(yy,\eta\eta)+p(yy)\mathbin{.}R(yy,\eta\eta)
\end{aligned}
\]
\Needspace{10\baselineskip}
\textit{it remains to note here (though the proof of this lies deeper)}
\[
\begin{aligned}&p(yy)\mathbin{.}\dfrac{\mathrm dr(yy)}{\mathrm d\log y}
-r(yy)\mathbin{.}\dfrac{\mathrm dp(yy)}{\mathrm d\log y}
\\[6pt]&=\tfrac12p(yy)\mathbin{.}r(yy)\mathbin{.}\{p(yy)^4-r(yy)^4\}\end{aligned}
\]
\Needspace{10\baselineskip}
\textit{hence}
\[
\begin{aligned}&\dfrac{P'P'-PP''}{PP}-\dfrac{S'S'-SS''}{SS}
\\[6pt]&=-\tfrac12\left(\dfrac{PP}{SS}+\dfrac{SS}{PP}\right)
\mathbin{.}p(yy)\mathbin{.}r(yy)\mathbin{.}\{p(yy)^2-r(yy)^2\}\end{aligned}
\]
\Needspace{14\baselineskip}
\textit{one also finds}
\[
\begin{aligned}
PS'-SP'={}&\tfrac12\left\{\begin{gathered}y^{\frac18}+3y^{\frac98}\\[3pt]-5y^{\frac{25}8}
 -7y^{\frac{49}8}+9y^{\frac{81}8}+\mathord{.}-\mathord{.}-\mathord{.}\end{gathered}\right\}\times\\[10pt]
&\times\left\{\begin{gathered}y^{\frac18}(\eta^{\frac12}+\eta^{-\frac12})
\\[3pt]-y^{\frac98}(\eta^{\frac32}+\eta^{-\frac32})
\\[3pt]-y^{\frac{25}8}(\eta^{\frac52}+\eta^{-\frac52})
 +\mathord{.}+\mathord{.}-\mathord{.}-\end{gathered}\right\}
\end{aligned}
\]
\Needspace{6\baselineskip}
\textit{The square of the second factor on the other side of the preceding equation becomes}
\[
=r\mathbin{.}Q(y,\eta\eta)+q\mathbin{.}R(y,\eta\eta)
\]
\Needspace{6\baselineskip}
\textit{The first factor becomes}
\[
=\{p(yy)^2+r(yy)^2\}{}^4\!\surd \tfrac12
 \{p(yy)^2-r(yy)^2\}p(yy)ryy)\ ?
\]
\Needspace{12\baselineskip}
\textit{Taken together, reductis reducendis, this becomes}


\SourceLeaf{411}

\[
\begin{aligned}
(PS'-SP')^2={}&\tfrac14\{p(yy)^2+r(yy)^2\}\\[6pt]&\quad\{p(yy)P(yy,\eta\eta)-r(yy)R(yy,\eta\eta)\}\times\\
&\times\{r(yy)P(yy,\eta\eta)+p(yy)R(yy,\eta\eta)\}\\[10pt]
={}&\tfrac14\{[p(yy)^2-r(yy)^2]PP+2p(yy)r(yy)SS\}\times\\
&\times\{2p(yy)r(yy)PP-[p(yy)^2-r(yy)^2]SS\}
\end{aligned}
\]
\Needspace{16\baselineskip}
\textit{If, therefore, one sets}
\[
\eta=e^{i\varphi}
\]
\[
\dfrac{-iS}{P}\surd {\dfrac{p(yy)^2-r(yy)^2}{2p(yy)r(yy)}}=\sin\theta
\]
\textit{this gives}
\[
\mathrm d\varphi=
\dfrac{2\,\mathrm d\theta}
 {\surd {\left(\begin{gathered}{}[p(yy)^2-r(yy)^2]\cos\theta^2\\[3pt]{}+
[p(yy)^2+r(yy)^2]\sin\theta^2\end{gathered}\right)}}
\]
\[
\cos\theta=\dfrac RP\surd {\dfrac{p(yy)^2+r(yy)^2}{2p(yy)r(yy)}}
\]
\begin{center}\rule{0.23\textwidth}{0.3pt}\end{center}
\clearpage
\endgroup
